\documentclass[bachelor, och, coursework]{SCWorks}
% Тип обучения (одно из значений):
%    bachelor   - бакалавриат (по умолчанию)
%    spec       - специальность
%    master     - магистратура
% Форма обучения (одно из значений):
%    och        - очное (по умолчанию)
%    zaoch      - заочное
% Тип работы (одно из значений):
%    coursework - курсовая работа (по умолчанию)
%    referat    - реферат
%  * otchet     - универсальный отчет
%  * nirjournal - журнал НИР
%  * digital    - итоговая работа для цифровой кафедры
%    diploma    - дипломная работа
%    pract      - отчет о научно-исследовательской работе
%    autoref    - автореферат выпускной работы
%    assignment - задание на выпускную квалификационную работу
%    review     - отзыв руководителя
%    critique   - рецензия на выпускную работу

% * Добавлены вручную. За вопросами к @mchernigin
\usepackage{preamble}

\begin{document}

% Кафедра (в родительном падеже)
\chair{математической кибернетики и компьютерных наук}

% Тема работы
\title{Разработка интерфейса образовательной платформы с использованием Svelte}

% Курс
\course{2}

% Группа
\group{251}

% Факультет (в родительном падеже) (по умолчанию "факультета КНиИТ")
% \department{факультета КНиИТ}

% Специальность/направление код - наименование
% \napravlenie{02.03.02 "--- Фундаментальная информатика и информационные технологии}
% \napravlenie{02.03.01 "--- Математическое обеспечение и администрирование информационных систем}
% \napravlenie{09.03.01 "--- Информатика и вычислительная техника}
\napravlenie{09.03.04 "--- Программная инженерия}
% \napravlenie{10.05.01 "--- Компьютерная безопасность}

% Для студентки. Для работы студента следующая команда не нужна.
% \studenttitle{студентки}

% Фамилия, имя, отчество в родительном падеже
\author{Толства Роберта Сергеевича}

% Заведующий кафедрой 
\chtitle{доцент, к.\,ф.-м.\,н.}
\chname{С.\,В.\,Миронов}

% Руководитель ДПП ПП для цифровой кафедры (перекрывает заведующего кафедры)
% \chpretitle{
%     заведующий кафедрой математических основ информатики и олимпиадного\\
%     программирования на базе МАОУ <<Ф"=Т лицей №1>>
% }
% \chtitle{г. Саратов, к.\,ф.-м.\,н., доцент}
% \chname{Кондратова\, Ю.\,Н.}

% Научный руководитель (для реферата преподаватель проверяющий работу)
\satitle{доцент, к.\,ф.-м.\,н.} %должность, степень, звание
\saname{С.\,В.\,Миронов}

% Руководитель практики от организации (руководитель для цифровой кафедры)
\patitle{доцент, к.\,ф.-м.\,н.}
\paname{С.\,В.\,Миронов}

% Руководитель НИР
\nirtitle{доцент, к.\,п.\,н.} % степень, звание
\nirname{В.\,А.\,Векслер}

% Семестр (только для практики, для остальных типов работ не используется)
\term{2}

% Наименование практики (только для практики, для остальных типов работ не
% используется)
\practtype{учебная}

% Продолжительность практики (количество недель) (только для практики, для
% остальных типов работ не используется)
\duration{2}

% Даты начала и окончания практики (только для практики, для остальных типов
% работ не используется)
\practStart{01.07.2024}
\practFinish{13.01.2024}

% Год выполнения отчета
\date{2025}

\maketitle

% Включение нумерации рисунков, формул и таблиц по разделам (по умолчанию -
% нумерация сквозная) (допускается оба вида нумерации)
\secNumbering

\tableofcontents

\defabbr
\begin{description}
  \item \textbf{SPA, Single Page Application} "--- одностраничное приложение
  \item \textbf{UI} "--- пользовательский интерфейс
  \item \textbf{UX} "--- пользовательский опыт
\end{description}

\intro

Современные образовательные платформы стали неотъемлемой частью цифровой трансформации образования, обеспечивая удобное взаимодействие студента и преподавателя через канал заданий, опросов и материалов \cite{intro_digital_transformation}. Однако эффективность таких платформ во многом зависит от качества их интерфейса: он должен быть интуитивным, отзывчивым и адаптивным для пользователей с разными потребностями. Также желательно, чтобы он в той или иной мере удовлетворял современным стандартам UI/UX-дизайна.

В условиях роста конкуренции на рынке EdTech становится необходимостью использовать передовые технологии, которые позволяют создавать высокопроизводительные и масштабируемые решения. Одной из таких технологий является Svelte "--- инновационный фреймворк, который предлагает принципиально новый подход к построению пользовательских интерфейсов за счет компиляции компонентов в чистый JavaScript на этапе сборки, что исключает необходимость виртуального DOM и сокращает накладные расходы.

Из всего вышеупомянутого можно выделить \textbf{актуальность темы}:

\begin{itemize}
  \item Рост спроса на онлайн-образование требует создания платформ с высокой производительностью и минимальной задержкой. Ситуация усугубляется общей низкой производительностью среднего SPA.
  \item Традиционные фреймворки, такие как React, Vue, Angular, сталкиваются с ограничениями в оптимизации рендеринга, что критично для реурсоемких приложений \cite{intro_compare_frameworks}.
  \item Svelte, благодаря своей архитектуре, позволяет достичь большей эффективности, что делает его перспективным инструментом для разработки высоконагруженных фронтендов (в том числе и для образовательных платформ) \cite{intro_svelte_docs}.
\end{itemize}

\textbf{Целью работы} является изучение технологии Svelte, а также разработка части графического интерфейса образовательной платформы MergeMinds с использованием этого фреймворка.

В ходе работы должны быть выполнены следующие \textbf{задачи}:
\begin{itemize}
  \item Провести анализ требований к интерфейсам образовательных платформ и существующих технологических решений.
  \item Изучить особенности Svelte, включая его экосистему и методы оптимизации.
  \item Спроектировать архитектуру интерфейса, реализовать ключевые модули (авторизация, курсовая навигация, интерактивные элементы)
  \item Создать взаимодействие с существующим backend-сервисом с использованием архитектуры REST API \cite{intro_rest_api}.
\end{itemize}

% После введения — серии \section, \subsection и т.д.

\conclusion


% Отобразить все источники. Даже те, на которые нет ссылок.
% \nocite{*}

\bibliographystyle{ugost2003}
\bibliography{thesis}

% Окончание основного документа и начало приложений Каждая последующая секция
% документа будет являться приложением
\appendix

\end{document}
